% 04_extension.tex - Phần 4: Triển khai Chrome Extension
\section{Triển khai Chrome Extension}

\subsection{Tổng quan kiến trúc}

Hệ thống Algorithmic Circuit Breaker v3.0 được triển khai dưới dạng một Chrome Extension độc lập kết hợp với Flask Backend siêu nhẹ quản lý mô hình kỹ thuật. Kiến trúc được thiết kế xoay quanh tính toán \textit{Shift-Left}:

\begin{figure}[H]
    \centering
    \begin{tikzpicture}[node distance=2.5cm, auto, thick, scale=0.85, transform shape]
        \tikzstyle{service} = [rectangle, draw, fill=green!20, text width=6em, text centered, minimum height=2.5em, rounded corners]
        \tikzstyle{db} = [cylinder, draw, fill=yellow!20, text width=5em, text centered, minimum height=2em, aspect=0.3]
        
        \node [service] (content) {content.js (NLP Local)};
        \node [service, right of=content, node distance=4cm] (background) {background.js};
        \node [service, right of=background, node distance=4cm] (flask) {Flask Server};
        \node [db, below of=flask, node distance=2.5cm] (state) {Session State};
        \node [service, above of=flask, node distance=2.5cm] (pid) {PID Controller};
        \node [service, left of=pid, node distance=4cm] (user) {RL UserAgent};
        
        \draw [<->] (content) -- node[above] {sendMessage} (background);
        \draw [<->] (background) -- node[above] {REST API} (flask);
        \draw [<->] (flask) -- (state);
        \draw [<->] (flask) -- (pid);
        \draw [<->] (pid) -- (user);
    \end{tikzpicture}
    \caption{Kiến trúc hệ thống Extension + Flask Backend (V3.0)}
    \label{fig:extension_arch}
\end{figure}

\subsubsection{Các thành phần chính}

\begin{table}[H]
\centering
\caption{Các thành phần của hệ thống V3.0}
\begin{tabular}{|l|p{5cm}|p{6cm}|}
\hline
Thành phần & Mô tả & Công nghệ \\
\hline
content.js & Service quản lý phân tích NLP trực tiếp trên DOM và thu thập Scrolling Velocity. Xử lý đồ họa hiển thị (Terminal UI). & JavaScript, CSS injection \\
\hline
background.js & Dịch vụ cầu nối trung gian (Service Worker) forward payload JSON bảo mật. & Chrome Extension API, fetch \\
\hline
server.py & Điểm cuối API (API Gateway). Liên kết Tác Tử Học Tăng Cường với Bộ Điều Khiển. & Python, Flask \\
\hline
manifest.json & Khung chỉ thị quyền hệ thống. & Manifest V3 \\
\hline
\end{tabular}
\end{table}

\subsection{Kiến trúc Client-Server Nguyên khối (Monolithic Client-Server Architecture)}

Trái ngược với các thiết kế phân tán cồng kềnh, nguyên mẫu (Proof-of-Concept) hiện tại được xây dựng hoàn toàn dựa trên khối lượng nhẹ (lightweight architecture). Chrome Extension vận hành như một "cảm biến mù" (blind sensor) ở tuyến đầu - chỉ thu thập vận tốc cuộn và điểm số độc hại vô hướng. Dữ liệu này được đóng gói thành một payload JSON cơ bản và truyền dẫn thông qua giao thức REST API đến một Backend Python Flask đơn nút mạng (single-node) tại \texttt{server.py}. Để tối ưu hóa khắt khe độ trễ sinh ra trong khối hệ thống điều khiển PID phản hồi nhanh, trạng thái phiên hoạt động của người dùng (user session state) được quản lý và xử lý trực tiếp trên bộ nhớ lưu trữ RAM dưới dạng cấu trúc từ điển nguyên sinh (in-memory dictionary), từ bỏ hoàn toàn việc sử dụng mạng lưới Vi dịch vụ (Microservices) hay các hệ quản trị CSDL trung gian phức tạp như Redis Cache. Kiến trúc này đặc biệt đáp ứng chuẩn mực vận hành thời gian thực.

\subsubsection{Mô hình UserAgent Học Tăng Cường}

Lớp UserAgent loại bỏ hoàn toàn suy đoán ngẫu nhiên thông qua hàm Rescorla-Wagner chính xác:

\begin{lstlisting}[language=Python, caption=Trích đoạn cấu trúc tác tử RL UserAgent V3.0]
class UserAgent:
    def __init__(self, config: RLConfig = None):
        self.config = config or RLConfig()
        self.expected_reward = self.config.initial_expected_reward
        self.rpe = 0.0
        
    def calculate_rpe(self, velocity: float, toxicity_score: float) -> float:
        norm_v = min(velocity / 800.0, 1.0)
        actual_reward = min(norm_v + (toxicity_score * 0.5), 1.0)
        
        # Rescorla-Wagner RPE Calculation
        self.rpe = actual_reward - self.expected_reward
        self.expected_reward = self.expected_reward + self.config.alpha * self.rpe
        return self.rpe
\end{lstlisting}

Điểm gây nghiện (Addiction Score) được chuẩn hóa thẳng từ RPE với phương pháp nhân hệ số giới hạn để đẩy vào mô hình PID. 

\subsection{Cấu trúc Chrome Extension} 

\subsubsection{Manifest V3 Configuration}

Để tuân thủ tuyệt đối quy định an ninh lưu trữ của Chrome Security:
\begin{lstlisting}[language=JavaScript, caption=manifest.json (Lược dịch)]
{
  "manifest_version": 3,
  "name": "Algorithmic Engagement Monitor",
  "version": "3.0.0",
  "permissions": ["activeTab", "scripting"],
  "host_permissions": ["http://127.0.0.1:5000/*"],
  "background": {
    "service_worker": "background.js"
  },
  "content_scripts": [{
    "matches": ["<all_urls>"],
    "js": ["content.js"],
    "run_at": "document_idle"
  }]
}
\end{lstlisting}

\subsubsection{Background Bridge Pattern}

Vấn đề: Do Content Scripts chạy trên bảo mật Content Security Policy (HTTPS, ví dụ: \textit{facebook.com}) không thể gửi gói trực tiếp tới máy chủ phát triển (HTTP Localhost). Nút trung chuyển được thiết lập để bẻ khóa chính sách CORS một cách an toàn thông qua \textit{Message Passing}.

\begin{lstlisting}[language=JavaScript, caption=Định tuyến Packet bảo mật]
chrome.runtime.onMessage.addListener((message, sender, sendResponse) => {
    if (message.type === 'SEND_METRICS') {
        const payload = message.payload;
        // Payload has strict { velocity, toxicity, scroll_time }
        fetch('http://127.0.0.1:5000/analyze', {
            method: 'POST',
            headers: { 'Content-Type': 'application/json' },
            body: JSON.stringify(payload)
        })
        .then(response => response.json())
        .then(data => sendResponse({ success: true, data: data }));
        return true; 
    }
});
\end{lstlisting}

Độ phức tạp quyền riêng tư được kiểm soát tối đa khi chỉ có số liệu trừu tượng $O(1)$ lọt vào đường ống Fetch. Mọi cụm từ khóa (keywords) trong DOM được loại khỏi giao diện gửi.

\subsubsection{Giao diện Terminal Monitor}

Thao tác UI (User Interface) được chuẩn hóa theo phong cách "Terminal/Monitor" loại bỏ các cụm từ y tế tiêu cực. Thông tin hệ thống chỉ có dạng các biến ma trận cảnh báo kỹ thuật số.

\begin{lstlisting}[language=CSS, caption=Trích dẫn giao diện Terminal UI]
.cb-gauge-val {
    position: absolute !important;
    bottom: 4px !important;
    left: 50% !important;
    transform: translateX(-50%) !important;
    font-size: 22px !important;
    font-weight: 700 !important;
    color: var(--cb-cyan) !important;
    text-shadow: 0 0 15px var(--cb-cyan) !important;
}
\end{lstlisting}

\subsection{Luồng dữ liệu (Data Pipeline) 3.0}

\begin{enumerate}
    \item Người dùng tương tác tốc độ cao trên DOM mạng xã hội.
    \item \textbf{content.js} giám sát vận tốc trung bình (2000ms moving window) và Quét Toxicity tại chỗ bằng tĩnh từ điển `TOXICITY_DICTIONARY`.
    \item Xây dựng Payload vô hình dạng Metadata (`{ velocity: 654, toxicity: 0.1, scroll_time: 2 }`).
    \item \textbf{background.js} truyền đi dưới chuẩn HTTP Fetch đến Cổng `/analyze`.
    \item \textbf{server.py} thực thi phương trình học tăng cường tính ra lượng dư RPE.
    \item Cung cấp thông số RPE, Toxicity, Thời gian Session vào PID Controller phân loại.
    \item Trả tín hiệu điều khiển trực tiếp về \textbf{content.js} để kích hoạt ngắt mạch DOM.
\end{enumerate}
